\part{Literary Form as an Experiment in Interpretation}

\chapter{Aphorism, Fragment, and the Interrupted Reader}

Nietzsche's literary forms belong to his philosophical method and cannot be treated as packaging for doctrines stated elsewhere. An aphorism can invite confident interpretation and then leave its grounds underdetermined. A sequence of sections can establish an expectation that a later fragment unsettles. A passage can be spoken in a voice whose relation to the author is left open. In each case the reader must decide whether to assimilate the disruption, to revise the interpretation under way, or to preserve incompatible possibilities, and what the reader does is the activity of interpretation that the earlier chapters described in the abstract. In the preface to the \emph{Genealogy} Nietzsche insists that an aphorism has not been deciphered merely because it has been read, and that its exegesis, which demands an art of its own, only then begins. The remark presents reading as work that the text solicits and does not complete.

Reader-response theory and philosophical hermeneutics supply resources for describing these activities. Iser's account of the act of reading emphasizes the gaps and indeterminacies that a text leaves for the reader to fill, and the way expectations formed at one point are modified by what follows \citep{iser1978}. Gadamer's hermeneutics treats understanding as the fusion of horizons in which the interpreter's prejudgments are brought into question by the text they bring to bear on it \citep{gadamer1989}. Nehamas has argued that Nietzsche's styles of writing are bound up with his conception of a life composed as a literary character, so that the form of the work enacts the view it expresses \citep{nehamas1985}. Predictive processing offers a restricted comparative vocabulary for the same phenomena, namely expectation, discrepancy and revision, and the aim of this chapter is to say what that vocabulary captures and what it risks flattening.

A minimal model of reading treats the reader as inferring, after each fragment, a posterior over interpretive hypotheses $h\in\mathcal{H}$. Let $f_{1:t}$ be the fragments read so far, and let the posterior be
\[
q_t(h)\;\propto\; p(h)\prod_{u\le t}p(f_u\mid h,f_{1:u-1}).
\]
An interruption is a fragment $f_{t+1}$ that is improbable under the currently dominant hypothesis but not under others, so that the posterior shifts. The three responses named above then have a formal counterpart. Assimilation retains the dominant hypothesis and absorbs the fragment by lowering its precision or redescribing it as a local exception. Revision transfers weight to a different hypothesis. Plural retention keeps several hypotheses alive, so that the entropy of $q_t$ is not allowed to collapse. The formalism does not, on its own, assign a value to this response. The entropy of $q_t$ is a posterior uncertainty about which interpretation obtains, which is a different quantity from the ambiguity term of Chapter~5, an expectation of the entropy of observations given states, and several incompatible readings can remain plausible even when each predicts the features of the text with precision. Neither posterior entropy nor likelihood ambiguity says what distinguishes productive interpretive plurality from uncertainty that contributes nothing to understanding, and literary judgment evidently draws that distinction. This is a first place where a generic discrepancy measure falls short of the practice it describes: it has no resources for saying which plurality counts as an achievement.

Irony exposes the limitation more sharply. An ironic utterance supports a content that is the opposite of its surface, and it does so through a stance that the reader must infer from cues. A faithful model therefore needs a latent variable for stance, $k\in\{\text{sincere},\text{ironic}\}$, so that hypotheses are pairs $h=(c,k)$ of a content and a stance, and the likelihood of the same words depends on both. The reader's inference about $k$ is itself learned from the author's habits, from genre and from the surrounding text. Once this variable is admitted, the single quantity called prediction error splits into discrepancies at different levels, and some of the most interesting effects in Nietzsche, in which a passage is at once endorsed and mocked, correspond to postures in which the reader cannot settle $k$ and in which the instability of $k$ is the point. Historical allusion and the double use of a term introduce further latents. Each such addition is a modelling decision that goes beyond the basic account, and one lesson of the chapter is that the adequacy of the predictive vocabulary for literature depends on how much structure is supplied to it by the analyst.

The material form of the text enters here, as it did in Chapter~1. A reader of the editorial arrangement of the notebooks meets sequences whose order is the editors', and may mistake the order for an argument. The consequence for interpretation is the same as that of a precision setting imposed from outside: the reader's expectations are shaped before reading begins, and the frame is mistaken for the content.

\chapter{Zarathustra's Voice and the Hazards of Discipleship}

\emph{Thus Spoke Zarathustra}, subtitled a book for all and none, multiplies the positions from which words are spoken. The chapter distinguishes Nietzsche as author, a narrator who reports Zarathustra's journeys and speeches, Zarathustra himself, the figures to whom he speaks, including his disciples, the animals who accompany him and the higher men who gather at his cave in the last part, and finally the reader, who may be addressed through any of these. It investigates how prophetic authority is constructed and then complicated by the narrative, and it declines to treat every declaration as a transparent statement of doctrine. A speech delivered by Zarathustra to a crowd that cannot understand him is a different kind of utterance from a speech delivered to his animals, and a speech whose terms the narrative subsequently undercuts has a different status again.

The central issue is the tension between teaching self-overcoming and creating followers attached to the teacher's authority. The text states the tension itself. In the prologue Zarathustra says that the creator seeks companions and not followers, and at the close of the speech on the bestowing virtue he asks his disciples to lose him and find themselves, promising to return only when all of them have denied him, and observing that a teacher is repaid poorly by a pupil who remains only a pupil. The device of departure and return is repeated across the work. Zarathustra withdraws to solitude, comes back to the company of others, finds that his teaching has been misunderstood or turned into doctrine, and leaves again. Later episodes continue the theme, among them the shadow who follows him, whom he tells to go free, and the gathering of the higher men, who in their relief at finding a master end by worshipping an ass. The comedy of the final scenes is evidence that the narrative itself treats the worship of the teacher as a hazard to be mocked and not as an outcome to be sought.

These features bear on the attachment-theoretical reading developed in Chapter~11. That reading holds, in the terms of the paper by Ferrari \citep{ferrari}, that Zarathustra's solitary self-overcoming resembles insecure-avoidant agency, with independence functioning as defensive self-sufficiency, little corrective feedback from relationships and insufficient acknowledgment of vulnerability. The text supplies evidence for such a reading. Zarathustra repeatedly retreats to solitude when his teaching is misunderstood, he is addressed in the work as one who longs for the solitude that he calls his home, and his friendships, where they appear, are frequently strained or marked by his impatience. A reader inclined to the attachment lens can accumulate such passages.

The text also supplies evidence that complicates it, and the interpretive test is whether the reading survives its own counterevidence. The speech on the bestowing virtue can be read as the rejection of dependent discipleship in favour of a relation between more independent participants, with the promised return indicating a relationship that continues across separation. Desire for relationship and defensive avoidance are not mutually exclusive, and the invitation to companions in the prologue or the repeated return to people therefore cannot by itself refute the avoidant reading, since an agent may long for company and still organize his conduct so as to limit what company can ask of him. The weight of the analysis falls on scenes of actual exchange, and these are mixed. In ``The Stillest Hour'' at the close of the second part, Zarathustra is addressed by a voice he cannot resist, confesses fear and says that he lacks the voice for commanding, and is told that he must still become a child, a scene of vulnerability and of being commanded, although the speaker is internal. In ``The Convalescent'' he collapses after confronting his most abysmal thought and lies for seven days while his animals attend to him and then speak with him, a scene of receiving care, though from animals and not from persons. His meetings with the wanderer who calls himself his shadow and with the higher men in the fourth part stage companionship and weariness together. None of this proves secure attachment, but it complicates the claim that separation expresses defensive withdrawal as a general matter. The textual basis the paper offers is also thin. Its characterization of Zarathustra as unilateral, isolated and unresponsive to relational feedback rests on two quotations, the proclamation of the overman in the prologue and the passage on burning in one's own flame, and neither is drawn from a scene in which Zarathustra and another person actually exchange anything. A further difficulty is that the paper's description of Zarathustra's agency as absolute, unmodulated and unresponsive to others fits poorly with scenes in which he is moved, exhausted or reproached by others.

One textual correction is needed. The paper locates the passage about burning in one's own flame in the speech titled ``On the Sublime Ones'' in the second part, whereas the passage belongs to ``On the Way of the Creating One'' in the first part. The correction is more than a matter of reference, because the image concerns transformation through self-overcoming, which is a statement about a process of becoming, and by itself supplies little evidence about whether a person can receive care or respond to relational feedback. Evidence of the second kind has to be found in scenes of actual exchange.

The conclusion of the chapter retains the question and weakens the verdict. The distinction that organizes the evidence is between independence that permits relationship and independence that must defend itself against relationship, and the narrative stages both without simply choosing between them. A reader who finds only avoidance has not accounted for the invitations and returns, and a reader who finds only openness has not accounted for the retreats and the contempt. The text's own refusal to settle the matter is consistent with the reading of Chapter~13, on which the form of the work asks the reader to hold competing interpretations of autonomy and relationship and to decide what to make of the instability.

\chapter{Metaphor, Concept Formation, and the Invention of Distinctions}

Metaphor is a possible means of reorganizing salience and conceptual structure, and Nietzsche gives it a foundational role. In the early essay ``On Truth and Lies in a Nonmoral Sense,'' which he did not publish \citep{nietzsche1979}, concepts are described as the residue of metaphors whose figurative origin has been forgotten, so that what passes for literal truth is the accumulated and worn-out product of transfers between spheres. Whatever one makes of the historical claim about language, the thought has a systematic consequence for the framework of this book. Schematization, in the sense of Chapter~2, requires a stock of distinctions, and the origin of new distinctions is the question that Chapter~9 left open. Metaphor is one candidate for where some of them come from.

Two functions can be distinguished. Some metaphors help a reader assimilate an unfamiliar subject to an existing category, so that the unfamiliar is handled by the repertoire already in place. Others make previously unnoticed relations available for thought, so that something is seen as having a structure it was not seen to have. Only the second function bears on structural revision. A way of stating the difference is that assimilating metaphors act within the current representation, whereas generative metaphors alter it.

The mathematical discussion returns to representation maps. Let $\phi:X\to Z$ be the existing organization and let a candidate representation be $\phi':X\to Z'$. Three relations between them correspond to three effects. If $\phi=g\circ\phi'$ for some map $g$, the new representation refines the old one and introduces distinctions that the old one lacked. If $\phi'=g\circ\phi$, it coarsens it and renders earlier distinctions dispensable. If neither holds, the two cut the state space along incomparable lines, and the new representation both preserves some distinctions and loses others. A metaphor can be modelled as a partial structure-preserving map $h$ from a source domain with relations $R_S$ to a target domain, proposing the relations $h(R_S)$ as new relations on the target. Whether the proposal is productive is then assessed by the criterion of Chapter~2: whether the representation it induces is sufficient, to within an acceptable margin, for the tasks that matter, and whether it opens tasks that the previous one could not formulate.

Model selection and minimum description length supply a second comparison. A description length criterion prefers the representation minimizing the sum of the length of the model and the length of the data given the model, in the manner of Rissanen's principle \citep{rissanen1978}. Such criteria capture a real virtue, the compression of regularities, and they might be expected to favour metaphors that unify disparate phenomena. Neither shorter descriptions nor greater predictive accuracy is assumed here to exhaust literary or philosophical value. A metaphor may lengthen the description of the data in the short term, since it introduces a structure whose payoff arrives only later or in another domain, and a metaphor's worth can lie in the questions it makes askable and not in the accuracy of its answers. The criteria also assume that the data to be described are given, whereas a generative metaphor can alter what counts as data.

The chapter therefore asks how a text can expand interpretive possibilities without immediately settling them. A metaphor that settles its relation to the target at once becomes a definition and loses its generative force, whereas one that leaves the transferred structure partly open invites the reader to continue the mapping. This returns the argument to the reader. The transformation of the reader's repertoire is accomplished through the reader's own work with the figure, and to that extent the text enacts the claim of Chapter~9 that new distinctions are not obtained by inference within a fixed family but made available to it. Whether such availability can be assessed by any standard other than the reader's subsequent capacities is a question that the final part carries into the discussion of shared models and institutions.
